\documentclass[letterpaper]{article}
\usepackage[submission]{aaai2027}
\usepackage[hyphens]{url}
\usepackage{graphicx}
\usepackage{listings}
\urlstyle{rm}
\def\UrlFont{\rm}
\usepackage{booktabs}
\usepackage{amsmath}
\frenchspacing

\setcounter{secnumdepth}{0}

\title{Supplementary Appendix for\\Temporal Hand Labanotation}
\author{Anonymous Submission}
\affiliations{}

\lstset{
    basicstyle=\ttfamily\small,
    breaklines=true,
    breakatwhitespace=false,
    columns=fullflexible,
    keepspaces=true
}

\newcommand{\ppath}[1]{\path{#1}}
\newcommand{\pathentry}[2]{\noindent\textbf{#1}\\\ppath{#2}\par}
\newcommand{\cmdentry}[2]{\noindent\textbf{#1}\\\texttt{python }\ppath{#2}\par}

\begin{document}
\maketitle

\section{Purpose}

This appendix records the concrete artifacts, protocols, and settings used by
the anonymous submission. It is intended as a reproducibility companion to the
main paper and a claim-to-evidence map for the reported results.

\section{Directory Layout}

All paths below are relative to the project root unless stated otherwise.

\pathentry{Submission package root:}{paper/submission_ready/}
\pathentry{Main paper file:}{paper/submission_ready/AnonymousSubmission2027.tex}
\pathentry{Checklist file:}{paper/submission_ready/ReproducibilityChecklist.tex}
\pathentry{Generated experiment artifacts:}{experiments/generated/}
\pathentry{Executable scripts:}{tools/*.py}
\pathentry{Base dataset root:}{/opt/tiger/InterHand}

\section{Claim-to-Evidence Map}

\textbf{Claim 1:} THL improves symbolic structure over a flat HL-style sequence
representation.
\begin{itemize}
    \item Main-paper evidence: Table 1A and the ablation table
    \item Source artifacts:
    \ppath{experiments/generated/topline_evidence_bundle.json}
    and
    \ppath{experiments/generated/temporal_representation_ablation_report.json}
\end{itemize}

\textbf{Claim 2:} THL yields substantially stronger local editability than
matched proxy token spaces.
\begin{itemize}
    \item Main-paper evidence: Table 1A
    \item Source artifacts:
    \ppath{experiments/generated/local_edit_audit.json},
    \ppath{experiments/generated/learned_token_proxy_report.json}
\end{itemize}

\textbf{Claim 3:} THL improves interaction-conditioned preservation under the
same held-out protocol.
\begin{itemize}
    \item Main-paper evidence: Table 1B and the statistical discussion in
    Results
    \item Source artifacts:
    \ppath{experiments/generated/topline_evidence_bundle.json},
    \ppath{experiments/generated/interaction_realized_significance_report.json},
    \ppath{experiments/generated/interaction_realized_feasible_left_repair_bundle.json}
\end{itemize}

\textbf{Claim 4:} Short transition chunks are the strongest minimal temporal
upgrade.
\begin{itemize}
    \item Main-paper evidence: Table 1C and the ablation table
    \item Source artifacts:
    \ppath{experiments/generated/interaction_realized_closing_chunk_transfer_val_only.json},
    \ppath{experiments/generated/interaction_realized_opening_chunk_transfer_val_only.json},
    \ppath{experiments/generated/temporal_representation_ablation_report.json}
\end{itemize}

The ablation table should be read within the learned sequence-model family
rather than as a direct replacement for every value in the main result table.
The main result table mixes structure, editability, and interaction-preservation
audits, whereas the ablation isolates which temporal channels matter inside the
matched sequence-model comparison.

\section{Temporal HL Export And Preprocessing}

The submission does not introduce a new raw image dataset. Instead, it reuses
the HL-derived symbolic export already built from InterHand2.6M and related
annotations in the current workspace.

\pathentry{Temporal export entry point:}{tools/build_temporal_hl.py}
\pathentry{Train split used in symbolic pretraining:}{experiments/generated/temporal_hl_train.json}
\pathentry{Validation split used for support construction and tuning:}{experiments/generated/temporal_hl_val.json}
\pathentry{Test split used for final evaluation:}{experiments/generated/temporal_hl_test.json}

The main preprocessing stages are:
\begin{itemize}
    \item conversion from joint annotations to symbolic hand states
    \item temporal windowing for learned sequence models
    \item support-bank construction for symbolic intervention and chunk-transfer audits
    \item slice filtering for hard interaction and feasible-subset evaluations
\end{itemize}

\section{Main Paper Numbers And Their Source Artifacts}

\pathentry{Topline evidence bundle used for Table 1:}{experiments/generated/topline_evidence_bundle.json}
\pathentry{Ablation report used for Table 2:}{experiments/generated/temporal_representation_ablation_report.json}
\pathentry{Locality audit:}{experiments/generated/local_edit_audit.json}
\pathentry{Interaction significance report:}{experiments/generated/interaction_realized_significance_report.json}
\pathentry{Feasible interaction repair bundle:}{experiments/generated/interaction_realized_feasible_left_repair_bundle.json}
\pathentry{Broader closing chunk-transfer report:}{experiments/generated/interaction_realized_closing_chunk_transfer_val_only.json}
\pathentry{Broader opening chunk-transfer report:}{experiments/generated/interaction_realized_opening_chunk_transfer_val_only.json}

\section{Main-Paper Values}

The paper reports the following principal values from the artifacts above.
\begin{itemize}
    \item structure: flat symbolic sequence $0.6923$ vs grouped symbolic sequence $0.8974$
    \item locality: symbolic locality $0.3553$ vs semantic proxy $0.0945$ and continuous proxy $0.0848$
    \item hard interaction, strict full slice: closing/opening $0.0250/0.0533$
    \item hard interaction, corrected support: closing/opening $0.1807/0.1812$
    \item feasible interaction subset: closing/opening $0.6392/0.6497$
    \item preserve-side repair: closing/opening $0.6582/0.6497$
    \item chunk transfer: closing/opening $0.7785/0.7197$ against fixed-edge $0.7025/0.6306$
\end{itemize}

\section{Protocol Fairness Notes}

The central comparisons in the main paper are designed to isolate the effect of
the representation rather than the effect of a stronger downstream solver. The
fairness conditions used throughout are:
\begin{itemize}
    \item the symbol inventory is unchanged when THL is compared against flat
    HL-style symbolic baselines
    \item the downstream operator family is held fixed within each audit
    (local intervention, interaction-conditioned preservation, or chunk
    transfer)
    \item support banks are built from validation data rather than test data
    \item final reporting is performed on held-out unseen-subject test splits
    \item learned-token controls use matched split families and stronger
    training budgets rather than weaker geometric heuristics
\end{itemize}

This means the reported gains should be interpreted as gains in notation
structure and editability under matched operator strength, not as gains from a
privileged retrieval or repair mechanism.

\section{Statistical Evidence}

The paper reports significance only for the interaction-conditioned frontier,
where frame-level variance is most consequential. The corresponding artifact is
\ppath{experiments/generated/interaction_realized_significance_report.json}.

The reported results are:
\begin{itemize}
    \item on the full hard slice, corrected support improves closing by
    $+0.156$ with 95\% CI $[0.127, 0.186]$ and opening by $+0.128$ with
    95\% CI $[0.101, 0.156]$; paired permutation tests are below $10^{-4}$ in
    both cases
    \item on the feasible interaction subset, preserve-side repair improves
    closing by $+0.184$ with 95\% CI $[0.108, 0.259]$ and opening by $+0.166$
    with 95\% CI $[0.108, 0.223]$; the corresponding $p$-values are
    $5\times10^{-5}$ and below $10^{-4}$
\end{itemize}

The sequence-accuracy comparisons additionally use three random seeds
$\{0,1,2\}$ for learned sequence models, while symbolic intervention and chunk
transfer are deterministic once the split and support bank are fixed.

\section{Claim Boundary}

The submission makes a deliberately narrow claim. It does not claim that THL is
a universal solution for every downstream hand-motion task, nor that the paper
exhausts every possible evaluation of symbolic hand notation. The supported
claim is instead:
\begin{itemize}
    \item HL can be upgraded from framewise notation into a temporal symbolic
    motion representation
    \item this upgrade improves structure, local editability, and
    interaction-conditioned preservation under matched held-out protocols
    \item the strongest minimal upgrade is grouped state plus short transition
    motifs, rather than heavier temporal channels
\end{itemize}

This is the scope the main paper is written to defend, and the artifacts listed
above are selected accordingly.

\section{Protocol Sketch}

For clarity, the overall evaluation protocol used in the paper can be summarized
as follows.
\begin{enumerate}
    \item Build temporal symbolic train/validation/test splits from the
    HL-derived export.
    \item Train the learned sequence models on the train split and select
    hyperparameters on validation.
    \item Build support banks only from validation data.
    \item Keep the downstream operator family fixed for the comparison being
    audited.
    \item Evaluate on the held-out test split, including unseen-subject slices
    and the interaction-conditioned subsets used in the main paper.
\end{enumerate}

This sketch is intentionally high level: the exact entry points and artifact
files are listed below, while the main paper focuses on the representation claim
rather than on engineering workflow detail.

\section{Common Reviewer Checks}

The present submission is designed to rule out three common alternative
interpretations.
\begin{itemize}
    \item \textbf{``This is just a stronger recognizer.''} The paper does not
    evaluate THL only through framewise prediction; it evaluates structure,
    locality, and interaction-conditioned preservation under fixed downstream
    operator classes.
    \item \textbf{``The gain comes from a larger symbol set.''} The grouped
    state component preserves the original HL symbol inventory; the measured
    gain comes from reorganizing the object and adding short transition motifs.
    \item \textbf{``The gain comes from a privileged solver.''} Support banks,
    operator families, and held-out split protocols are matched across the
    comparisons used for the main claims.
\end{itemize}

\section{Final Hyperparameters Reported In The Main Paper}

\textbf{Learned temporal symbolic sequence models.}
\begin{itemize}
    \item seeds: $\{0,1,2\}$
    \item window size: $32$
    \item stride: $16$
    \item hidden dimension: $128$
    \item pretraining epochs: $20$
    \item finetuning epochs: $20$
\end{itemize}

\textbf{Matched learned-token controls.}
\begin{itemize}
    \item compared token codebooks: $32$ and $64$
    \item final reported codebook size: $64$
    \item max codebook frames: $100{,}000$
    \item pretraining window span: $186$
    \item pretraining step: $96$
    \item pretraining epochs: $120$
    \item finetuning epochs: $120$
\end{itemize}

\textbf{Chunk transfer.}
\begin{itemize}
    \item tested chunk lengths: $2, 3, 4$
    \item broader opening result is stable across all three lengths
    \item broader closing result is reported at the strongest tested setting
\end{itemize}

\section{Representative Commands}

\textbf{Matched old-HL vs temporal-HL boundary check.}
\cmdentry{Command:}{tools/build_oldhl_temporal_matched_report.py}

\textbf{Grouped/repair structural analysis.}
\cmdentry{Command 1:}{tools/build_current_code_symbolic_frontier.py}
\cmdentry{Command 2:}{tools/build_representation_risk_bundle.py}

\textbf{Broader feasible closing chunk transfer.}
\cmdentry{Command:}{tools/build_interaction_realized_closing_chunk_transfer.py}

\textbf{Broader feasible opening chunk transfer.}
\cmdentry{Command:}{tools/build_interaction_realized_opening_chunk_transfer_val_only.py}

\textbf{Sequence-model training family.}
\cmdentry{Command 1:}{tools/train_symbolic_pretrain.py}
\cmdentry{Command 2:}{tools/train_symbolic_pretrain_grouped.py}
\cmdentry{Command 3:}{tools/train_joint_token_baseline.py}

\section{Script Inventory Used Directly By The Submission}

The submission relies primarily on the following executable entry points.

\pathentry{Entry point 1:}{tools/build_temporal_hl.py}
\pathentry{Entry point 2:}{tools/build_current_code_symbolic_frontier.py}
\pathentry{Entry point 3:}{tools/build_representation_risk_bundle.py}
\pathentry{Entry point 4:}{tools/build_topline_evidence_bundle.py}
\pathentry{Entry point 5:}{tools/build_oldhl_temporal_matched_report.py}
\pathentry{Entry point 6:}{tools/build_interaction_realized_closing_chunk_transfer.py}
\pathentry{Entry point 7:}{tools/build_interaction_realized_opening_chunk_transfer_val_only.py}
\pathentry{Entry point 8:}{tools/train_symbolic_pretrain.py}
\pathentry{Entry point 9:}{tools/train_symbolic_pretrain_grouped.py}
\pathentry{Entry point 10:}{tools/train_joint_token_baseline.py}
\pathentry{Entry point 11:}{tools/build_paper_figures.py}

\section{Selection Criteria}

The final settings used in the main paper were not chosen post hoc from all
available experiments. We used the following rule:
\begin{itemize}
    \item for matched learned controls, keep the setting with the strongest
    validation-side sequence accuracy under the same split protocol
    \item for chunk transfer, keep the setting with the strongest corrected
    feasible-slice score, but only if the effect remains stable across the
    tested chunk lengths or support-bank variants
    \item keep only evidence that directly supports the paper's central
    representation claim under the held-out protocol
\end{itemize}

\section{Environment}

The current environment used to generate the final paper bundle is:
\begin{itemize}
    \item Python $3.11.2$
    \item PyTorch $2.9.1$
    \item CUDA-enabled execution path in the current training scripts
    \item NVIDIA A100 80GB PCIe GPUs in the active machine environment
\end{itemize}

\section{Anonymity And Scope}

The submission package contains only anonymous paper assets, bibliography
files, figures, checklist material, and this appendix. Internal rendering
snapshots, exploratory debug images, and unrelated workspace artifacts are not
part of the final submission bundle. The present appendix documents only the
artifacts required to understand or reproduce the reported anonymous results.

\end{document}
